% 00-declaration.tex
\addchap{Selbstständigkeitserklärung}

Wir versichern hiermit wahrheitsgemäß, dass wir das vorliegende \DocName{} mit dem Titel
\begin{quote}
  \textbf{\Title} \\
  \ifdefstring{\Subtitle}{NoSubtitle}{}{%
    \textit{\Subtitle}
  }
\end{quote}
selbstständig und ohne unzulässige fremde Hilfe verfasst und keine anderen als die angegebenen Quellen und Hilfsmittel benutzt haben. Alle Stellen, die wörtlich oder sinngemäß aus veröffentlichten oder nicht veröffentlichten Schriften entnommen sind, haben wir als solche kenntlich gemacht.

Die Arbeit ist in gleicher oder ähnlicher Form noch bei keiner anderen Prüfungsbehörde eingereicht und auch nicht veröffentlicht worden.

\section*{Erklärung zum Einsatz generativer KI-Werkzeuge}

Diese Arbeit wurde unter Einsatz generativer Text- und Code-KI-Werkzeuge erstellt. Durch diese Erklärung am Anfang der Arbeit gelten das Dokument und die nachfolgend aufgeführten Punkte als gekennzeichnet; auf eine kleinteilige Markierung einzelner Sätze im Haupttext wurde zugunsten des Leseflusses verzichtet.

\textbf{Verwendete Werkzeuge:} Claude (Anthropic), Gemini (Google) und ChatGPT (OpenAI).

\begin{itemize}
  \item \textbf{Quellcode (Kapitel~\ref{chap:hauptteil}):} Die Simulationsumgebung, die Auswertungsskripte und die Notebook-Zellen wurden unter Einsatz generativer Coding-Werkzeuge erstellt. Sämtliche Bestandteile wurden vom Projektteam geprüft, ausgeführt und fachlich verantwortet; die berichteten Messwerte stammen ausnahmslos aus eigenen Läufen.
  \item \textbf{Formatierung und Struktur (global):} Die Werkzeuge wurden zur Unterstützung der \LaTeX-Formatierung, der Erstellung von Tabellenstrukturen sowie zur Aufbereitung von Quellcode-Darstellungen eingesetzt.
  \item \textbf{Sprachliche Überarbeitung (global):} Im gesamten Dokument wurden die Werkzeuge unterstützend für Rechtschreibprüfung, Grammatikkorrektur und Formulierungsvorschläge eingesetzt.
  \item \textbf{Formulierungshilfe:} Die Kurzfassung, der englische Abstract, die Zusammenfassung sowie die Ausformulierung der in den Projektnotizen stichpunktartig festgehaltenen Ergebnisse wurden unter Einbezug KI-basierter Formulierungsvorschläge erarbeitet und anschließend vom Projektteam fachlich geprüft.
\end{itemize}

\vspace{2.5cm}

\noindent
\Location, den \DeliveryDate

\vspace{2.0cm}

\noindent
\begin{tabular}{@{}p{6.5cm}@{\hspace{1.5cm}}p{6.5cm}@{}}
  \hrulefill & \hrulefill \\
  \NameA     & \NameB     \\
\end{tabular}

\clearpage
